\section{Conclusion}

This paper presented a staged TextVQA study that was designed to be broader than a single tuned-model report while remaining methodologically disciplined under a realistic compute budget. The zero-shot funnel established Qwen2.5-VL-3B as the strongest open-source backbone in the benchmark pool, the finalist reruns showed that this conclusion survived transfer to the official validation split, and the winner-conditioned LoRA study converted that backbone advantage into a tuned validation score of 84.76\%. The resulting paired gain over the strongest zero-shot finalist is 6.04 points on the same 5,000 official validation questions, with a bootstrap 95\% confidence interval of $[5.00, 7.04]$ points. The main empirical conclusion is therefore anchored on official validation evidence, while the internal-dev split is used for screening and diagnostics.

The empirical findings are also more specific than a single leaderboard improvement. Prompting matters, but mostly within a backbone rather than across backbones. Fine-tuning helps, but it helps most on OCR-grounded questions rather than on OCR-free ones. Training-stage loss is a useful signal for allocating follow-up budget, but not a substitute for answer-accuracy-based promotion or official validation evaluation. OCR toggling matters less after adaptation than it does in zero-shot prompting, while additional training data appears to help modestly in the diagnostic follow-up runs. Together, these observations make the final tuned model more interpretable than a generic ``best checkpoint'' claim.

From a project-design perspective, the work also illustrates a practical template for academic multimodal studies. The key ingredients are not exotic: a stratified internal-dev split with an explicit split audit, a clear promotion rule, a hard boundary between local evaluation and remote training, a winner-conditioned follow-up phase, and artifact-level reproducibility. Yet in combination, these decisions make the final evidence substantially easier to defend. Future extensions should add a semantically richer learned judge with human calibration, expand adaptation beyond a single backbone, and test whether the same funnel protocol transfers to adjacent OCR-heavy benchmarks. Even without those extensions, the current study supports a clear conclusion: carefully adapted open-source VLMs can become materially stronger TextVQA systems when the experimental protocol separates screening evidence from final held-out validation evidence.
